\documentclass[11pt]{article}
\usepackage[margin=1in]{geometry}
\usepackage{amsmath}
\usepackage{booktabs}
\usepackage{url}
\setlength{\emergencystretch}{2em}
\title{A finite-encounter terminal set with lane-hold backup controllers}
\date{Prepared October 8, 2026}
\begin{document}
\maketitle

\section{Request}
The user asked to explore a terminal set for a finite encounter: the safety
task ends when the target leaves the perception range, modelled as an
absorbing mode of a hybrid system, and the terminal set is the set of states
from which a backup controller completes the encounter safely, so that the
prediction horizon need not cover the whole avoidance. The proposal also
suggested a persistent-separation set from a collision-cone condition, and
asked whether the backup should only accept already separating states or
also continue avoiding. The user asked for an implementation and a
simulation test if the design is reasonable.

\section{Assessment}
\begin{itemize}
\item The CLF-tube terminal set of October 6 already is a safe-exit set: its
  backup is the CLF terminal controller (return to the given path at the
  reference speed), and the tube must miss the target until it leaves the
  range. What the proposal adds is the finite-encounter formulation without
  the 60-s sliding window (hypothesis H4), and a backup that may keep
  avoiding.
\item \emph{Finite encounter.} Requiring the target to leave the range
  removes H4, and $S\cup\{\dagger\}$ is invariant with an absorbing end state.
  It also makes every state non-terminal for a target that stays in range: a
  parallel target at 49.9~m (two unit tests) and the first frames of the
  braking lead, whose first estimate is a lead at the ego's speed ($A=0$ until
  0.5~s). A fixed encounter window keeps both properties: the encounter ends at
  the earlier of the target's exit and a fixed time 60~s after the encounter's
  start. Because that end does not move, the shift argument holds within the
  encounter without H4; a target still present then starts a new encounter.
\item \emph{Backup controller.} The return-to-path backup makes no state
  terminal while a slower vehicle stays in the given lane; this is the
  braking lead's startup failure in all 12 noisy runs. A family of lane-hold
  backups (the same CLF and tube about each lane centre of the road) lets the
  terminal set include overtaking states. The tube certificate selects a backup
  whose tube stays clear, which is the role the proposal gave the collision
  cone.
\item \emph{Persistent separation.} $\mathcal S_{\mathrm{sep}}$ is the
  separating terminal set of October 4--6, withdrawn because it is not
  invariant under the ego's return to the path. It was not implemented; the
  tube check covers both cases the proposal distinguished.
\item \emph{Return to nominal.} A lane backup kept alone would never be left:
  the shifted endpoint keeps its lane and appended holds keep it there (the
  conflict that withdrew the shoulder terminal on October 4). A plan ending
  in another lane is therefore extended by path guidance to the nominal set
  whenever that extension meets every hard row.
\end{itemize}
The lanes are modes taken from the road description; this relaxes the
October 6 requirement that the terminal set use no mode.

\section{Implementation}
\begin{itemize}
\item \path{controller/terminalSafeSet.m}: backups (\texttt{modeReferences},
  one reference per lane centre of \texttt{road.laneOffsets}, nominal first;
  on a curve the offset path's trim with the given path's CLF matrix), the
  tube about a lane centre (station rate (3.2')), membership in the union
  (\texttt{member}, nominal first, cheap skip of backups whose level is above
  their region), \texttt{select}, the fixed window (\texttt{windowSteps};
  \texttt{clearToWindowEnd}, \texttt{encounterWindowEnded}), and the backup's
  holds.
\item \path{controller/nonlinearBicycleModel.m}: the CLF error and the path
  guidance subtract the reference's \texttt{lateralOffset}.
\item \path{controller/solvePredictiveControl.m}: the terminal cone uses the
  selected backup, a shifted plan is extended by holds of its endpoint's
  backup, and \texttt{localReturnToNominal} extends a plan that ends in
  another lane back to the nominal set when every appended hold meets the
  hard rows; the search records the anchor length, the return and the
  terminal lane.
\item \path{controller/collisionAvoidanceController.m}: the encounter start
  in the controller state and the remaining window in the model.
\item \path{controller/readControllerInputs.m} reads an optional
  \texttt{road.laneOffsets}, and the campaign road of
  \path{scripts/collisionThreatScenario.m} has four lane centres.
\item Theory: \path{controller/TERMINAL_SAFE_SET.md} (Sections 1--7 and 10).
\end{itemize}

\section{Tests}
\texttt{terminalSafeSetTest} now checks the fixed window (a parallel target is
cleared to the window end, a node after it is past the encounter, a short
remaining window ends the encounter before a slower lead is reached), that a
lane-hold backup makes terminal an overtaking state the nominal backup cannot
complete, the invariance of that backup over 40 holds, and the offset CLF
error. Before the fixed window, two tests (a parallel target at 49.9~m in
\texttt{nonlinearPredictiveSafetyTest} and the window test) failed with the
exit-only set. The full suite result is in \path{tests.txt}.

\section{Campaigns}
The 14 encounters with exact states and the 84 noisy ones (seeds 20261003 and
1--5), 2000-hold cap, 1-s recovery dwell, each in its own MATLAB process:
\begin{itemize}
\item \emph{baseline}: \texttt{5b4379a} controller (identical to
  \texttt{68883d2});
\item \emph{A}: exit required within 60~s of each node, with lane backups;
\item \emph{D}: fixed window with lane backups (the committed design);
\item \emph{E}: fixed window without lane backups.
\end{itemize}
\begin{center}
\begin{tabular}{lrrrr}
\toprule
 & baseline & A & D & E\\
\midrule
exact, recovered of 14 & 14 & 14 & 14 & 14\\
noisy, recovered of 84 & 55 & 56 & 57 & 55\\
noisy, new recoveries / new failures & -- & 2 / 1 & 3 / 1 & 0 / 0\\
collisions & 0 & 0 & 0 & 0\\
noisy smallest clearance (m) & 0.272 & 0.289 & 0.289 & 0.272\\
noisy runs using a lane backup & -- & 4 & 13 & 0\\
\bottomrule
\end{tabular}
\end{center}
Per run: \path{TERMINAL_BACKUP_SET_20261008/campaign-runs.csv}; summary
\path{comparison.txt}.

\paragraph{Fixed window.} E reproduces the baseline run for run: the window
end is never reached within these encounters, so the window removes H4
without changing behaviour.

\paragraph{Lane backups.} With exact states the braking lead uses a lane
backup in about 80 of its 250 frames (the overtaking) and recovers, as
before. With noise, the braking lead failed at its first frame in all 12
baseline runs or at the next (solver time limit: no endpoint in the terminal
set; at most 1 hold). With D, 9 of 12 run beyond 1 hold (9 to 1865 holds) and
one recovers. The
other stops are numerical (\texttt{pcbfNoNumericalResult},
\texttt{clfNoNumericalResult}) or \texttt{nonlinearAcceptanceFailed} later in
the run. Two inspected frames show the causes: a 479-hold startup plan whose
endpoint drifts by about 2~m when rolled out from the next estimate, followed
by a stalled CLF stage (Clarabel status 9, residual $3\times10^{-5}$); and a
plan that swerves to 9.7~m left, whose primal problems are then infeasible at
every trust scale. With A (exit required) the braking lead cannot start at
all. The new failure (8~m/s curved head-on, seed 1, in A and D) occurs at hold
1571, long after the encounter; E recovers it at hold 1764.

\paragraph{Conclusion.} The finite-encounter terminal set with lane-hold
backups is sound and implemented, and the fixed window removes H4 at no cost.
The lane backups make the braking lead's startup feasible, but the noisy
campaign gains only two recoveries (55 to 57): its dominant failures are not
an empty terminal set but forecast changes on long open-loop plans and
numerical stalls of the conic solver.

\section{Decision needed}
\begin{itemize}
\item Keep the lane backups, or the nominal backup alone. The lane backups
  are modes, against the October 6 requirement; the return to the nominal
  path is a guidance extension, not a certificate. (The window question is
  settled by the addendum: no preset time is needed.)
\item The remaining failures point to the plan's length and the conic
  solver's numerical stalls, not to the terminal set.
\item Whether to add a startup rollout toward the other side (braking lead at
  startup, addendum).
\end{itemize}

\section{Reproduction}
Scripts in \path{TERMINAL_BACKUP_SET_20261008/scripts}: \texttt{campaign.sh}
(with \texttt{SEEDS}, \texttt{D}, \texttt{VARIANT}), \texttt{run-A.sh},
\texttt{run-DE.sh}, \texttt{runOneEncounter.m} (seed 0 runs with exact
states), \texttt{compare.py}, and \texttt{diagnoseFrame*.m} with
\texttt{debugControllerModel.m} (the controller's model construction) for
failing frames. Each variant ran on a snapshot of the source: A is the
implementation before the fixed window (the target had to leave the range
within 60~s of each node), D the source of commit \texttt{b836fba}, E that
source without \texttt{laneOffsets} in the scenario road, and F the source
with proven separation (the addendum). F ran as D in \texttt{run-DE.sh}
(\texttt{D} the study folder, \texttt{VARIANT='struct()'}, \texttt{SEEDS}
unset, i.e.\ seeds 0, 20261003 and 1--5) with the F snapshot as source; the
summary is \texttt{compare.py} over \texttt{out-base out-A out-D out-E out-F}. Their hashes are in
\path{variant*-sha256.txt}. Raw outputs were in
\path{~/.cache/collisionAvoidance/terminal_backup_20261008/} and are
disposable.

\section{Addendum: proven separation instead of a preset window}
The user asked why the window is 60~s and whether a preset time is needed.
It is not: 60~s was the engineering cap of October 6, with no derivation, and
the safety task itself has no duration. The encounter now ends at the first of
two computed events, the target's exit or a proven separation, and
\texttt{terminal.horizonSeconds} only bounds how far the check computes; a
state that reaches it with neither event is not terminal (variant F, which
replaces D in the source).

\paragraph{Proven separation.} On a straight road the distance between two
path boxes is at least their gap along the road and their gap across it.
After a grid time the ego band across the road only shrinks, and along the
road it moves at the trim's station rate and widens by at most
$T\sigma(r)$ more. A straight-line target ($\beta=0$) keeps a gap across the
road open if its lateral velocity never closes it ($sV\ge0$ and $sA\ge0$ on
the side away from the ego, and opposite signs on the other) and a gap along it
open if its station rate never closes it ($cV\ge v$ and $cA\ge0$ ahead,
$cV\le v$ and $cA\le0$ behind). A circling target stays in the disk of its
circle widened by its body reach. These hold for all later times, also
through a stop. A separation proven for a state holds for its successors
under the backup at the same absolute time, so invariance and the shift
argument hold within an encounter without any window. On a curve only the
exit ends an encounter.

\paragraph{Implementation.} \path{controller/terminalSafeSet.m}:
\texttt{permanent} (the certificate above, tested at every grid point of a
chunk), \texttt{localTargetMotion} (the forecast's tangential acceleration and,
for $\beta\ne0$, the circle's centre and widened radius), the target's course
and speed on the grid, and the computation bound \texttt{horizonSteps}; the
check ends at the first exit or proven separation after confirming the grid
points before it and the separation's own point, and returns
\texttt{noExitOrSeparation} when it reaches the bound with neither. The fixed
window (\texttt{windowSteps}, \texttt{clearToWindowEnd},
\texttt{encounterWindowEnded}) and the encounter start in the controller state
are removed. The issued plan and the trace record how the encounter ends
(\texttt{terminalEnd}). \texttt{terminalSafeSetTest} now checks that a
parallel target, a faster lead and a target circling beside the road end the
encounter by proven separation, that a slower lead in the ego lane meets the
tube, and that a lead closing only after the 60-s bound is not terminal. The
full suite (592 tests) passes on F (\path{tests.txt}).

\begin{center}
\begin{tabular}{lrrrrr}
\toprule
 & baseline & A & D & E & F\\
\midrule
exact, recovered of 14 & 14 & 14 & 14 & 14 & 14\\
noisy, recovered of 84 & 55 & 56 & 57 & 55 & 58\\
noisy, new recoveries / new failures & -- & 2 / 1 & 3 / 1 & 0 / 0 & 4 / 1\\
noisy braking leads recovered of 12 & 0 & 0 & 1 & 0 & 2\\
collisions & 0 & 0 & 0 & 0 & 0\\
\bottomrule
\end{tabular}
\end{center}
F's new recoveries are the braking leads of seeds 3 and 5 at 8~m/s, the
15~m/s head-on of seed 4 and the 8~m/s curved crossing of seed 4; its new
failure is the same late stop as in A and D (8~m/s curved head-on of seed 1,
at hold 1571). Over all 98 runs, of F's accepted endpoints with a target
8144 end the encounter by the target's exit and 2689 by proven separation; the
end lies at most 26.2~s after the endpoint (median 0, 99th percentile 8.7~s;
exact runs at most 7.5~s), so the 60-s computation bound was never reached. Five of the 12 noisy braking leads run beyond 1 hold with
F (9 with D). Per run: \path{campaign-runs.csv} (all five variants); summary
\path{comparison-all.txt}.

\paragraph{Braking lead at startup.} F stops 6 of the 12 noisy braking leads
at the first frame, D 2. The four that D started (seeds 1 and 20261003 at
both speeds; D failed them later, after 1 to 924 holds) have one cause. The
first estimate is a lead 6.2~m ahead, 0.03 to 0.10~m/s slower than the ego,
with a heading of 0.6 to 3.8~mrad toward the left. A settled ego in the given
lane meets it. In the right lane separation is proven at once: the lead's
lateral velocity points away. In the two left lanes the forecast lead creeps
toward the tube: it meets the first one's tube in two runs, and otherwise no
separation can be proven, while at that speed difference the lead does not
leave the range within the 60~s computed. The startup potential-field rollout
passes on the left (to 5.6--8.4~m), reaches no terminal state in its 512 holds,
and no plan is accepted; D, which accepted a tube clear until the window end,
started all four. The cause is the side of the startup rollout, not an empty
terminal set. A second startup rollout toward the other side, when the first
reaches no terminal state, is the obvious remedy; it is not implemented
(\path{brakingLead-startup-F.txt}; \path{scripts/whyNotTerminal.m},
\path{scripts/whySeedFails.m} with the copy written by
\path{scripts/makeSolveDiag.py}).

\end{document}
